% =====================================================================
%  ch00a.tex —— 绪论（原书 p.1–8）  Lead 自译
% =====================================================================

\chapter*{绪论}
\addcontentsline{toc}{chapter}{绪论}
\markboth{绪论}{}

% [原文 p.1]
代数几何从诞生之日起，人们就明白：双有理等价的簇共享许多性质。因此，在每一个双有理等价类中找出某种意义下最简单的簇，再仔细研究这些簇，是再自然不过的事。

每条不可约曲线都双有理等价于唯一一条非奇异射影曲线，因此研究非奇异射影曲线，无非就是研究所有曲线的双有理等价类。

曲面情形更为复杂：每个不可约曲面都双有理等价于无穷多个非奇异射影曲面。20 世纪初意大利代数几何学派发展起来的曲面极小模型理论，其目标是：从每个双有理等价类中选出唯一一个非奇异射影曲面。具体做法极其简单：若非奇异射影曲面含有一条自交数为 $-1$ 的非奇异有理曲线，则可将该曲线收缩成一点，得到一个新的非奇异射影曲面。把这一操作尽可能多地重复，多数情形下就得到唯一的“极小模型”。偶尔会得到不唯一的模型，但这些情形反而可以非常明确地描述出来。
% [Tutor Note: 这里的“收缩成一点”是 Castelnuovo 判据的反向使用：$(-1)$ 曲线可由 Castelnuovo 定理收缩，且收缩后仍是非奇异射影曲面，这正是曲面理论比高维理论幸运的地方。]

把这一方法推广到高维的尝试开始得晚得多。一个原因大概是：某些例子提示，类似的做法在高维行不通。

1980 年代初，M. Reid 与森重文指出：如果允许处理的对象不只是非奇异簇，还包括带有某种温和奇点的簇，那么高维极小模型理论或许可行。这套方法称为极小模型理论，也称森纲领。经过 X. Benveniste、川又雄二郎、J. Kollár、Reid、V. V. Shokurov、角田秀一、E. Viehweg、宫冈洋一等许多人的贡献之后，1988 年森重文完成了三维情形。

% [原文 p.2]
此后，这一纲领发展成一套能够成功应用于代数几何中许多问题的方法。

本书的目标，是帮助读者入门极小模型理论的技巧与思想。

第 1 章从有理曲线的存在定理出发，用几何的方法引入整个纲领。后者的结果在后面用得不多，但作为概念上的基础是重要的。

第 2 章仍是引论性质，讨论奇点的几个侧面，以及小平消没定理的一些有用的推广。

纲领的第一个主要部分——锥定理，在第 3 章证明。这个证明在所有维数都有效。

本书余下的部分实质上用于研究三维翻转与三维翻模。翻转与翻模是只在三维才首次出现的新型双有理变换。曲面与三维簇之间的大部分重大差别，都可以追溯到翻转与翻模。
% [Tutor Note: “只在三维才首次出现”要这样理解：二维时小收缩就是收缩 $(-1)$ 曲线或 $(-2)$ 曲线（后者是 Du Val 奇点的收缩），情形已经封闭；三维才第一次出现收缩的例外集是曲线、且该曲线不能整条收缩的“翻转”现象。]

第 4 章讨论某一类曲面奇点的分类。这些结果在研究三维理论时要用到。

极小模型纲领进行途中出现的奇点，在第 5 章研究。这些结果在任意维数都成立。

翻模在第 6 章研究。翻模比翻转容易理解，至少在三维，它的描述是令人满意的。

第 7 章讨论三维翻转。一般的三维翻转仍然过于复杂、证明也太长，不适于写进教科书。因此本书把讨论限制在一个特殊的类——半稳定翻转。在这种情形下，证明可以大幅简化。半稳定翻转在许多场合自然出现，对若干应用来说这已经够用。

各章内容的更详细说明，放在各章的开头。

§4.5 与 §5.5 都偏离了论证的主线。这两种情形下，我们都觉得只靠引用文献来交待所需结果是不够的，而由本书作些补充，对读者恐怕更好。

% [原文 p.3]
\section*{必要的预备知识}
\addcontentsline{toc}{section}{必要的预备知识}

本书假定读者具备 [Har77] 程度的基本代数几何知识。不过，正如下文所述，书中会假定一些结果，所以只要暂且接受 Zariski 主定理、以及仅在 §4.3 用到的“上同调与基变换定理”等极少几个结论，那么仅凭参考文献中列举的日文教科书的知识，也完全读得下去。§5.1 等处用到的环论结果，在 [永田74]、[松村80] 等标准教科书中（一部分在 [宫西90] 中）有证明。

本书有几处不加证明地使用了某些结果。

定理 1.10 的证明需要射的形变空间维数的估计公式。这个结果的证明是技术性的，但对 §1.1 与 §1.2 的许多论证本质上是必要的。这些定理在后面的节里并不需要。

§1.5 复习除子交点数的一些基本性质，并说明后面频繁使用的弱形式的 Riemann–Roch 定理。

§2.4 中叙述并使用代数上同调与解析上同调的基本比较定理，以及奇异上同调的 Hodge 分解的一个特殊情形。

所有这些情形，我们只引用结论，并不需要其证明中用到的技术。

有时会用到 Leray 谱序列（关于其证明，参见 [God58, 4.17]、[石井97, §2.6]，或 [宫西90, 2.28]）。% [Tutor Note: 英文原著此处并列三处出处：[God58, 4.17]、[HS71, VIII.9]、[Bre97, IV.6]；日译本以后两种换成本国的 [石井97, §2.6] 与 [宫西90, 2.28]。故 [HS71]、[Bre97] 在本书正文中不再出现，但仍保留在书末参考文献表中。]

\begin{theorem*}{0.1}
$f\colon X\to Y$ 是概形的射，$F$ 是 $X$ 上的拟凝聚层。此时存在如下的谱序列：
\[
E_2^{i,j}=H^i(Y,R^jf_*F)\ \Longrightarrow\ H^{i+j}(X,F).
\]
\end{theorem*}

本书在许多地方还用 [Hir64] 的奇点解消。这个结果以如下两种形式使用。（参见 [石井97, §3.3]、[川又97, §2.3]。）

% [原文 p.4]
\begin{theorem*}{0.2}[广中定理]
$X$ 是 $\C$ 上的代数簇（或是解析空间的紧子集的足够小的邻域），$I\subset\cO_X$ 是定义闭子概形（或子簇）$Z$ 的凝聚理想层。此时存在非奇异簇（或解析空间）$Y$ 以及射影态射 $f\colon Y\to X$，满足以下条件：
\begin{enumerate}[label=(\arabic*),leftmargin=2.4em,itemsep=1pt,topsep=2pt]
\item $f$ 在 $X\setminus\bigl(\Sing(X)\cup\Supp Z\bigr)$ 上是同构；
\item $f^*I\subset\cO_Y$ 等于可逆层 $\cO_Y(-D)$；
\item $\Exc(f)\cup D$ 是单正规交叉除子。
\end{enumerate}
\end{theorem*}

由 [Hir64] 的主定理 I、II（解析情形为主定理 I$'$、II$'$）可得。

把上述定理中 (1) 的部分去掉所得的结果，称为\textbf{广中弱定理}（weak Hironaka theorem），而本书需要的正是这个较弱的结果。关于拟射影簇 $X$ 的广中弱定理，[AdJ97]、[BP96]、[Par98] 中有非常短的证明。这些论文都把广中弱定理归结为 [KKMSD73] 的环面嵌入理论（torus embedding theory）。（这些论文只说明 $D$ 是单交叉除子，但其论证可用于证明主张 (3)。）[石井97, §3.4]、[小田85, 第 1、2 章] 中有相关的解说。
% [Tutor Note: “弱”只弱在放弃 (1)——即允许在 $\Sing(X)\cup\Supp Z$ 上也作改动。本书大多数论证只用到这弱形式，所以可以把广中定理当作黑箱。]

奇点解消的相对版本如下。

\begin{theorem*}{0.3}[广中定理]
$f\colon X\to C$ 是从 $\C$ 上的代数簇（或是解析空间的紧子集的足够小的邻域）到非奇异曲线的平坦（flat）态射，$B\subset X$ 是除子。此时存在从非奇异簇 $Y$ 出发的射影双有理态射 $g\colon Y\to X$，使得对所有的 $c\in C$，$\Exc(g)+g^*B+(f\circ g)^*(c)$ 都是单正规交叉除子。
\end{theorem*}

这可由 [Hir64] 的主定理 II（或 II$'$）推出。这个结果只在第 7 章 $\dim X=3$ 的情形用到。若 $C$ 与 $X$ 都是射影的，它是 [AK97, Thm. 2.1]（广中弱定理）的特殊情形。后一论文同样忽略了 $\Exc(g)$，但只要稍加修改其证明，就能在计入 $\Exc(g)$ 的情形下也奏效。

\begin{notation}{0.4}
为避免误解，下面列举本书使用的标准记号。
\end{notation}

% [原文 p.5]
\begin{enumerate}[label=(\arabic*),leftmargin=2.6em,itemsep=2pt,topsep=2pt]
\item[(0)] 说（代数或解析）簇（variety）时，通常假定它既约（irreducible）且被约（reduced）。说概形（scheme）或解析空间（analytic space）时，不作这些假定。但分离性（separatedness）总假定成立。

\item $X$ 是正规概形。素除子（prime divisor）是既约且被约的余维数 $1$ 的子概形。$X$ 上的除子（divisor）是 $X$ 上素除子的形式线性组合 $D=\sum d_iD_i$（其中 $d_i\in\Z$）。使用这种写法时，假定 $D_i$ 互不相同。$\Q$-除子（$\Q$-divisor）是素除子的形式线性组合 $D=\sum d_iD_i$（其中 $d_i\in\Q$）。若对所有的 $i$ 都有 $d_i\ge 0$，则称 $D$ 是有效（effective）的。对 $\Q$-除子 $A,B$，若 $A-B$ 有效，则记作 $A\ge B$ 或 $B\le A$。（由于这种记法容易与“$A-B$ 是 nef”的主张混淆，故不常用。）除子（或 $\Q$-除子）$D$ 若对某个 $0\ne m\in\Z$ 有 $mD$ 是 Cartier 的，则称为 $\Q$-Cartier 的。簇 $X$ 若所有 $\Q$-除子都是 $\Q$-Cartier 的，则称为 $\Q$-分解的（$\Q$-factorial）。除子 $D=\sum d_iD_i$ 的支集（support）是闭子概形 $\bigcup_{d_i\ne 0}D_i$，记作 $\Supp D$。

\item 两个除子 $D_1,D_2$ 线性等价（linear equivalence）时，记作 $D_1\sim D_2$。两个 $\Q$-除子 $D_1,D_2$ 数值等价（numerical equivalence）时，记作 $D_1\equiv D_2$。（$\Q$-除子的线性等价性不作定义。）$\Q$-除子 $D$ 若 $D\sim 0$（或 $D\equiv 0$），则称为平凡（trivial）（或数值平凡（numerically trivial））。

\item 固有概形上的 $\Q$-Cartier 除子 $D$，若对任意不可约曲线 $C$ 都有 $(D\cdot C)\ge 0$，则称为 nef 的。

\item 概形之间的态射（morphism）处处有定义，用实线箭头 $f\colon X\to Y$ 表示。概形之间的有理映射（rational map，或简称 map）只在稠密开集上有定义，用虚线箭头 $f\colon X\dashrightarrow Y$ 表示。

% [原文 p.6]
\item $f\colon X\to Y$ 是射，$D_1,D_2$ 是 $X$ 上的两个除子。若存在 $Y$ 上的 Cartier 除子 $B$ 使 $D_1\sim D_2+f^*B$，则称二者线性 $f$-等价（linearly $f$-equivalent），记作 $D_1\sim_f D_2$。若 $D_1,D_2$ 是 $\Q$-除子，且存在 $Y$ 上的 $\Q$-Cartier $\Q$-除子 $B$ 使 $D_1\equiv D_2+f^*B$，则称二者数值 $f$-等价（numerically $f$-equivalent），记作 $D_1\equiv_f D_2$。若 $D\sim_f 0$（或 $D\equiv_f 0$），则称 $D$ 是（线性）$f$-平凡的（或数值 $f$-平凡的（numerically $f$-trivial））。

\item 对概形 $X$，$\mathrm{red}\,X$ 表示支集与 $X$ 一致的那个唯一的约化闭子概形。

\item 对双有理态射 $f\colon X\to Y$，例外集（exceptional set）$\Exc(f)\subset X$ 是使 $f$ 不成为双正则的点 $\{x\in X\}$（即 $f^{-1}$ 在 $f(x)$ 处不是射的点 $x$）的全体。通常给 $\Exc(f)$ 带上约化结构，视作子概形。

\item $X$ 是非奇异簇，$D=\sum d_iD_i$ 是 $X$ 上的 $\Q$-除子。若各 $D_i$ 都非奇异、且它们处处横截地（transversally）相交，则称 $D$ 是单正规交叉（simple normal crossing）除子（简称 snc 除子）。若 $X$ 的每一点 $x$ 都存在足够小的欧氏邻域 $U\ni x$ 使 $D\cap U$ 是 snc 除子，则称 $D$ 是正规交叉（normal crossing）除子（简称 nc 除子）。

\item $X$ 是概形。$X$ 的全体奇点构成的子集记作 $\Sing(X)$。$X$ 的奇点解消（resolution）是从非奇异 $Y$ 出发的固有双有理态射 $g\colon Y\to X$。

\item $X$ 是概形，$D=\sum d_iD_i$ 是 $X$ 上的 $\Q$-除子。$(X,D)$ 的对数奇点解消（log resolution）是固有双有理态射 $g\colon Y\to X$，其中 $Y$ 非奇异、$\Exc(g)$ 是除子、且 $\Exc(g)\cup g^{-1}(\Supp D)$ 是 snc 除子。由定理 0.2，特征 $0$ 的域上的簇存在对数奇点解消。

% [原文 p.7]
\item $f\colon X\dashrightarrow Y$ 是概形的双有理映射。$Z\subset X$ 是使 $f$ 在稠密开集 $Z^0\subset Z$ 上双正则的闭子概形。$f(Z^0)$ 的闭包称为 $Z$ 的\textbf{双有理变换}（birational transform）。（有时也称为固有（proper）变换或严格（strict）变换。）记作 $f_*(Z)$。若 $g\colon Y\to X$ 是双有理态射，则对 $g^{-1}\colon X\dashrightarrow Y$ 要用与通常记法略有不同的 $g_*^{-1}(Z)$。进而，若 $X,Y$ 正规且 $Z$ 是不可约除子，则当 $\dim f(Z)<\dim Z$ 时约定 $f_*Z=0$，由此对 $\Q$-除子 $Z$ 定义 $\Q$-除子 $f_*Z$。

\item 对实数 $d$，其下取整（round down）是不超过 $d$ 的最大整数，记作 $\lfloor d\rfloor$；上取整（round up）是不小于 $d$ 的最小整数，记作 $\lceil d\rceil$；小数部分（fractional part）是 $d-\lfloor d\rfloor$，常记作 $\{d\}$。若 $D=\sum d_iD_i$（$D_i$ 为互不相同的素除子）是实系数除子，则定义 $D$ 的下取整为 $\lfloor D\rfloor:=\sum\lfloor d_i\rfloor D_i$，上取整为 $\lceil D\rceil:=\sum\lceil d_i\rceil D_i$，小数部分为 $\{D\}:=\sum\{d_i\}D_i$。

\item 若 $X$ 是解析空间，通常取任意紧子集 $Z\subset X$，在其足够小的开邻域 $U\supset Z$ 上展开论证。若论证上方便，也会不加说明地取更小的 $U$。$U$ 称为 $X$ 在 $Z$ 处的芽（germ）。若 $g\colon Y\to X$ 是解析空间之间的固有态射，通常也在上述 $U\subset X$ 上展开论证。在这种设定下，用于代数簇的论证往往照搬可用，上面引入的记法也同样使用。有理型映射（meromorphic map）与双有理型映射（bimeromorphic map）就简称为有理映射与双有理映射。

\item $P:=R$，或 $R=:P$，表示新记号 $P$ 定义为等于旧记号 $R$。

\item $\Z_{>0}$ 表示正整数全体（不含 $0$）。同样，$\R_{\ge 0}$ 表示非负实数全体。

% [原文 p.8]
\item 本书中，证明的结尾记作 $\blacksquare$，定义、例等的结尾记作 $\square$。又，定理、命题等，在不给出证明时以 $\blacksquare$ 结尾，在证明于后文给出时以 $\square$ 结尾。
% [Tutor Note: 这条与传统中文教材的用法相反，务必记牢：定理末尾的实心方块表示“本书不证”，空心方块表示“证明在后面”。下文所有陈述环境的尾符都照原书，不作改动。]
\end{enumerate}
